% ============================================================
% Question Bank: ch08-03 Comparing Two Populations Visually (Modified — FL/IN/MN/VA)
% Topic Title: Comparing Two Populations Visually (Multiple Display Emphasis)
% CCSS: 7.SP.B.3 (FL/IN/MN/VA: multiple display emphasis)
% Questions: 27 (15 MC + 6 SA + 6 Graphical)
% ============================================================

% --- Multiple Choice Questions (15) ---

\begin{practiceQuestion}{m-8-3-q01}{mc}
\begin{questionText}
Which type of display keeps all individual data values and allows comparison of two groups side by side?
\end{questionText}
\choiceA{Histogram}
\choiceB{Circle graph}
\choiceC{Back-to-back stem-and-leaf plot}
\choiceD{Frequency table}
\correctAnswer{C}
\explanation{A back-to-back stem-and-leaf plot shows all individual values for two data sets sharing the same stems, allowing direct comparison.}
\end{practiceQuestion}

\begin{practiceQuestion}{m-8-3-q02}{mc}
\begin{questionText}
Two dot plots show test scores for Class A and Class B. Class A's dots are clustered around $78$ and Class B's dots are clustered around $85$. What can you conclude?
\end{questionText}
\choiceA{Class A has higher scores on average.}
\choiceB{Class B has higher scores on average.}
\choiceC{Both classes have the same average.}
\choiceD{You cannot compare the two classes.}
\correctAnswer{B}
\explanation{The center of Class B's distribution ($85$) is higher than Class A's ($78$), indicating Class B has higher scores on average.}
\end{practiceQuestion}

\begin{practiceQuestion}{m-8-3-q03}{mc}
\begin{questionText}
A double bar graph compares test scores for two schools. School X has more students scoring $90$--$100$ and fewer scoring $60$--$69$ than School Y. What does this suggest?
\end{questionText}
\choiceA{School Y performed better overall.}
\choiceB{School X performed better overall.}
\choiceC{Both schools performed the same.}
\choiceD{The data cannot be compared.}
\correctAnswer{B}
\explanation{Having more students in higher score ranges and fewer in lower ranges suggests School X performed better overall.}
\end{practiceQuestion}

\begin{practiceQuestion}{m-8-3-q04}{mc}
\begin{questionText}
Two box plots are drawn on the same number line. Plot A has a longer box than Plot B. What does this indicate?
\end{questionText}
\choiceA{Plot A has a higher median.}
\choiceB{Plot A has more data points.}
\choiceC{Plot A has a greater interquartile range.}
\choiceD{Plot A has a smaller range.}
\correctAnswer{C}
\explanation{The length of the box represents the IQR (Q3 $-$ Q1). A longer box means a greater IQR, indicating more spread in the middle $50\%$ of the data.}
\end{practiceQuestion}

\begin{practiceQuestion}{m-8-3-q05}{mc}
\begin{questionText}
A circle graph shows that Team A spends $40\%$ of practice on drills and $30\%$ on scrimmages. Team B spends $25\%$ on drills and $50\%$ on scrimmages. Which statement is supported by the data?
\end{questionText}
\choiceA{Team A spends more time on scrimmages.}
\choiceB{Team B spends a greater proportion of time on drills.}
\choiceC{Team A allocates a larger share of practice to drills.}
\choiceD{Both teams spend the same proportion on drills.}
\correctAnswer{C}
\explanation{Team A: $40\%$ drills. Team B: $25\%$ drills. Team A allocates a larger share to drills.}
\end{practiceQuestion}

\begin{practiceQuestion}{m-8-3-q06}{mc}
\begin{questionText}
Two histograms show the heights of two basketball teams. Team 1 is roughly symmetric and Team 2 is skewed right. Which team is more likely to have a few very tall players pulling the distribution?
\end{questionText}
\choiceA{Team 1}
\choiceB{Team 2}
\choiceC{Both equally}
\choiceD{Cannot determine}
\correctAnswer{B}
\explanation{A right-skewed distribution has a tail stretching toward higher values, indicating a few very tall players are pulling the data in that direction.}
\end{practiceQuestion}

\begin{practiceQuestion}{m-8-3-q07}{mc}
\begin{questionText}
A frequency table shows that Group A has $15$ values in the $70$--$79$ range and Group B has $8$ values in that range. If both groups have $30$ total values, which group has a higher relative frequency for the $70$--$79$ range?
\end{questionText}
\choiceA{Group A}
\choiceB{Group B}
\choiceC{Both the same}
\choiceD{Cannot determine}
\correctAnswer{A}
\explanation{Group A: $\frac{15}{30} = 50\%$. Group B: $\frac{8}{30} \approx 26.7\%$. Group A has the higher relative frequency.}
\end{practiceQuestion}

\begin{practiceQuestion}{m-8-3-q08}{mc}
\begin{questionText}
When comparing two populations using box plots, the overlap of the boxes indicates:
\end{questionText}
\choiceA{the populations are identical}
\choiceB{there is no difference between the populations}
\choiceC{some data values from both groups fall in the same range}
\choiceD{the data is invalid}
\correctAnswer{C}
\explanation{Overlapping boxes show that the middle $50\%$ of data from each group shares some of the same values. It does not mean the groups are identical — it indicates similarity in that region.}
\end{practiceQuestion}

\begin{practiceQuestion}{m-8-3-q09}{mc}
\begin{questionText}
Two dot plots compare ages of volunteers at two shelters. Shelter A: ages spread from $14$ to $62$. Shelter B: ages spread from $25$ to $40$. Which shelter has greater variability?
\end{questionText}
\choiceA{Shelter A}
\choiceB{Shelter B}
\choiceC{Both the same}
\choiceD{Cannot tell from dot plots}
\correctAnswer{A}
\explanation{Shelter A's range is $62 - 14 = 48$ years, while Shelter B's range is $40 - 25 = 15$ years. Greater spread means greater variability.}
\end{practiceQuestion}

\begin{practiceQuestion}{m-8-3-q10}{mc}
\begin{questionText}
A back-to-back stem-and-leaf plot compares two data sets. The stems are shared in the middle. If one side has leaves $\{2, 5, 7\}$ on stem $4$ and the other side has leaves $\{1, 3, 8, 9\}$ on stem $4$, what are the data values on the second side?
\end{questionText}
\choiceA{$4, 3, 8, 9$}
\choiceB{$14, 13, 18, 19$}
\choiceC{$41, 43, 48, 49$}
\choiceD{$4.1, 4.3, 4.8, 4.9$}
\correctAnswer{C}
\explanation{In a stem-and-leaf plot, combine the stem ($4$) with each leaf: $41, 43, 48, 49$.}
\end{practiceQuestion}

\begin{practiceQuestion}{m-8-3-q11}{mc}
\begin{questionText}
Two histograms both show $25$ data values. Histogram A has most data between $50$ and $70$. Histogram B has data evenly spread from $30$ to $90$. Which histogram shows data with less variability?
\end{questionText}
\choiceA{Histogram A}
\choiceB{Histogram B}
\choiceC{Both have the same variability.}
\choiceD{Cannot determine from histograms.}
\correctAnswer{A}
\explanation{Histogram A's data is concentrated in a smaller range ($50$--$70$), showing less variability. Histogram B's data is spread across a wider range ($30$--$90$), showing greater variability.}
\end{practiceQuestion}

\begin{practiceQuestion}{m-8-3-q12}{mc}
\begin{questionText}
A circle graph compares how two cities allocate their budgets. City A: Education $45\%$, Public Safety $30\%$, Other $25\%$. City B: Education $30\%$, Public Safety $40\%$, Other $30\%$. Which city devotes a larger proportion to public safety?
\end{questionText}
\choiceA{City A}
\choiceB{City B}
\choiceC{Both the same}
\choiceD{Cannot determine}
\correctAnswer{B}
\explanation{City B allocates $40\%$ to public safety, compared to City A's $30\%$. City B devotes a larger proportion.}
\end{practiceQuestion}

\begin{practiceQuestion}{m-8-3-q13}{mc}
\begin{questionText}
Two box plots show delivery times. Company A: median $= 3$ days, IQR $= 2$ days. Company B: median $= 5$ days, IQR $= 1$ day. Which statement is accurate?
\end{questionText}
\choiceA{Company A is faster and more consistent.}
\choiceB{Company A is faster but less consistent.}
\choiceC{Company B is faster and more consistent.}
\choiceD{Company B is faster but less consistent.}
\correctAnswer{B}
\explanation{Company A has a lower median ($3$ days), so it is faster. But its IQR ($2$ days) is larger than Company B's ($1$ day), meaning delivery times vary more — it is less consistent.}
\end{practiceQuestion}

\begin{practiceQuestion}{m-8-3-q14}{mc}
\begin{questionText}
Two frequency tables show quiz scores. Class P has frequencies: $70$--$79$: $4$, $80$--$89$: $10$, $90$--$99$: $6$. Class Q has: $70$--$79$: $8$, $80$--$89$: $7$, $90$--$99$: $5$. Which class has a higher percentage of students scoring $80$ or above?
\end{questionText}
\choiceA{Class P}
\choiceB{Class Q}
\choiceC{Both the same}
\choiceD{Cannot determine}
\correctAnswer{A}
\explanation{Class P: $\frac{10+6}{20} = \frac{16}{20} = 80\%$. Class Q: $\frac{7+5}{20} = \frac{12}{20} = 60\%$. Class P has a higher percentage scoring $80$ or above.}
\end{practiceQuestion}

\begin{practiceQuestion}{m-8-3-q15}{mc}
\begin{questionText}
When visually comparing two populations, the gap between the medians is most meaningful when compared to the:
\end{questionText}
\choiceA{sum of the medians}
\choiceB{measure of variability such as the MAD or IQR}
\choiceC{total number of data points}
\choiceD{range of the larger data set}
\correctAnswer{B}
\explanation{The difference in medians should be measured relative to the variability (MAD or IQR). A gap that is large compared to the variability indicates a meaningful difference.}
\end{practiceQuestion}

% --- Short Answer Questions (6) ---

\begin{practiceQuestion}{m-8-3-q16}{sa}
\begin{questionText}
Two data sets have medians of $40$ and $52$. The MAD of each data set is approximately $4$. Express the difference in medians as a multiple of the MAD and state whether the difference is meaningful.
\end{questionText}
\correctAnswer{The difference is $12$, which is $\frac{12}{4} = 3$ MADs. This is a meaningful difference since $3$ MADs is large.}
\explanation{Difference: $52 - 40 = 12$. Divide by MAD: $12 \div 4 = 3$ MADs. A difference of $2$ or more MADs generally indicates a meaningful difference between populations.}
\end{practiceQuestion}

\begin{practiceQuestion}{m-8-3-q17}{sa}
\begin{questionText}
A back-to-back stem-and-leaf plot shows: Group X leaves $\{8, 5, 3\}$ and Group Y leaves $\{2, 4, 7, 9\}$ on stem $6$. What is the median of Group Y's values on this stem?
\end{questionText}
\correctAnswer{$65.5$}
\explanation{Group Y values on stem $6$: $62, 64, 67, 69$. With $4$ values, the median is the average of the middle two: $\frac{64 + 67}{2} = 65.5$.}
\end{practiceQuestion}

\begin{practiceQuestion}{m-8-3-q18}{sa}
\begin{questionText}
Two box plots show finishing times (in seconds) for two relay teams. Team A: Min $= 48$, Q1 $= 52$, Median $= 55$, Q3 $= 59$, Max $= 64$. Team B: Min $= 50$, Q1 $= 54$, Median $= 58$, Q3 $= 60$, Max $= 63$. Which team has faster typical times, and which team is more consistent?
\end{questionText}
\correctAnswer{Team A has faster typical times (lower median $55$ vs $58$). Team B is more consistent (smaller IQR: $60 - 54 = 6$ vs $59 - 52 = 7$).}
\explanation{Team A median: $55$ s (faster). Team A IQR: $59 - 52 = 7$. Team B IQR: $60 - 54 = 6$. Team B has a smaller IQR, so it is more consistent.}
\end{practiceQuestion}

\begin{practiceQuestion}{m-8-3-q19}{sa}
\begin{questionText}
Two histograms both show $20$ data points. Histogram A shows a bell-shaped (symmetric) distribution centered at $75$. Histogram B shows a uniform distribution from $60$ to $90$. Which data set has less variability? Explain.
\end{questionText}
\correctAnswer{Histogram A has less variability. A bell-shaped distribution has most values concentrated near the center ($75$), while a uniform distribution spreads values evenly across the entire range.}
\explanation{In a symmetric bell shape, data clusters around the mean, giving a smaller spread. A uniform distribution has equal frequency everywhere, using the full range, and thus higher variability.}
\end{practiceQuestion}

\begin{practiceQuestion}{m-8-3-q20}{sa}
\begin{questionText}
A circle graph for School A shows $55\%$ of students participate in after-school activities. A circle graph for School B shows $40\%$. School A has $400$ students and School B has $500$. Which school has more students in after-school activities?
\end{questionText}
\correctAnswer{School A: $0.55 \times 400 = 220$ students. School B: $0.40 \times 500 = 200$ students. School A has more.}
\explanation{Even though School B is larger, School A has a higher participation rate. $220 > 200$, so School A has more students in after-school activities.}
\end{practiceQuestion}

\begin{practiceQuestion}{m-8-3-q21}{sa}
\begin{questionText}
Two frequency tables show daily step counts for two fitness classes. Class 1 data is mostly in the $8{,}000$--$9{,}999$ range. Class 2 data is mostly in the $5{,}000$--$6{,}999$ range. If both classes have the same number of participants, which class is generally more active? How can you tell from the frequency table?
\end{questionText}
\correctAnswer{Class 1 is more active. The intervals with the highest frequencies are in higher step-count ranges ($8{,}000$--$9{,}999$) compared to Class 2 ($5{,}000$--$6{,}999$).}
\explanation{The interval with the greatest frequency indicates where most data is concentrated. Class 1's peak in a higher range means its members typically take more steps.}
\end{practiceQuestion}

% --- Graphical Questions (3 GMC + 3 GSA) ---

\begin{practiceQuestion}{m-8-3-q22}{gmc}
\begin{questionText}
The dot plots below show the number of push-ups completed by students in two PE classes.

\begin{center}
\begin{tikzpicture}
  % Class 1
  \node[above] at (4.5,1.8) {\textbf{Class 1}};
  \draw[->] (0,0) -- (10.5,0);
  \foreach \x in {0,1,...,10} {\draw (\x,0.1) -- (\x,-0.1) node[below]{\tiny \x};}
  % Data: 3,4,4,5,5,5,6,6,6,6,7,7,7,8,8
  \foreach \x/\y in {3/1,4/1,4/2,5/1,5/2,5/3,6/1,6/2,6/3,6/4,7/1,7/2,7/3,8/1,8/2}
    {\fill[funBlue] (\x,\y*0.3+0.15) circle (0.1);}
\end{tikzpicture}
\end{center}
\vspace{0.3cm}
\begin{center}
\begin{tikzpicture}
  % Class 2
  \node[above] at (4.5,1.8) {\textbf{Class 2}};
  \draw[->] (0,0) -- (10.5,0);
  \foreach \x in {0,1,...,10} {\draw (\x,0.1) -- (\x,-0.1) node[below]{\tiny \x};}
  % Data: 1,2,2,3,3,3,4,4,4,4,5,5,5,6,6
  \foreach \x/\y in {1/1,2/1,2/2,3/1,3/2,3/3,4/1,4/2,4/3,4/4,5/1,5/2,5/3,6/1,6/2}
    {\fill[funRed] (\x,\y*0.3+0.15) circle (0.1);}
\end{tikzpicture}
\end{center}

What is the approximate difference in the medians of the two classes?
\end{questionText}
\choiceA{$1$}
\choiceB{$2$}
\choiceC{$3$}
\choiceD{$4$}
\correctAnswer{B}
\explanation{Class 1 has $15$ values. Median = 8th value $= 6$. Class 2 has $15$ values. Median = 8th value $= 4$. Difference: $6 - 4 = 2$.}
\end{practiceQuestion}

\begin{practiceQuestion}{m-8-3-q23}{gmc}
\begin{questionText}
The box plots below compare the weights (in pounds) of two dog breeds.

\begin{center}
\begin{tikzpicture}
  % Breed A: min=20, Q1=30, med=40, Q3=50, max=65
  \draw[->] (0,0) -- (8.5,0);
  \foreach \x/\v in {0/10,1/20,2/30,3/40,4/50,5/60,6/70,7/80}
    {\draw (\x,0.1) -- (\x,-0.1) node[below]{\tiny \v};}
  \node[left] at (-0.3,1) {\small Breed A};
  \draw[thick] (1,0.7) -- (1,1.3); % min=20
  \draw[thick] (1,1) -- (2,1);     % whisker
  \draw[thick,fill=funBlue!30] (2,0.7) rectangle (4,1.3); % Q1=30 to Q3=50
  \draw[thick,funRed] (3,0.7) -- (3,1.3); % med=40
  \draw[thick] (4,1) -- (5.5,1);   % whisker
  \draw[thick] (5.5,0.7) -- (5.5,1.3); % max=65

  % Breed B: min=35, Q1=45, med=55, Q3=60, max=75
  \node[left] at (-0.3,2.2) {\small Breed B};
  \draw[thick] (2.5,1.9) -- (2.5,2.5); % min=35
  \draw[thick] (2.5,2.2) -- (3.5,2.2); % whisker
  \draw[thick,fill=funGreen!30] (3.5,1.9) rectangle (5,2.5); % Q1=45 to Q3=60
  \draw[thick,funRed] (4.5,1.9) -- (4.5,2.5); % med=55
  \draw[thick] (5,2.2) -- (6.5,2.2); % whisker
  \draw[thick] (6.5,1.9) -- (6.5,2.5); % max=75
\end{tikzpicture}
\end{center}

Which breed has a greater IQR?
\end{questionText}
\choiceA{Breed A}
\choiceB{Breed B}
\choiceC{They have the same IQR.}
\choiceD{Cannot determine from box plots.}
\correctAnswer{A}
\explanation{Breed A IQR: $50 - 30 = 20$ lbs. Breed B IQR: $60 - 45 = 15$ lbs. Breed A has the greater IQR ($20 > 15$).}
\end{practiceQuestion}

\begin{practiceQuestion}{m-8-3-q24}{gmc}
\begin{questionText}
The back-to-back stem-and-leaf plot below compares scores for two groups.

\begin{center}
\begin{tabular}{r|c|l}
\textbf{Group X} & \textbf{Stem} & \textbf{Group Y} \\
\hline
$8\;5\;2$ & $5$ & $3\;6\;9$ \\
$9\;6\;3\;1$ & $6$ & $0\;4\;7$ \\
$7\;4$ & $7$ & $1\;5\;8\;9$ \\
$5$ & $8$ & $2\;6$ \\
\end{tabular}
\end{center}

Which group has the higher median?
\end{questionText}
\choiceA{Group X}
\choiceB{Group Y}
\choiceC{They have the same median.}
\choiceD{Cannot determine.}
\correctAnswer{B}
\explanation{Group X values (sorted): $52, 55, 58, 61, 63, 66, 69, 74, 77, 85$. Median (10 values): $\frac{63 + 66}{2} = 64.5$. Group Y values (sorted): $53, 56, 59, 60, 64, 67, 71, 75, 78, 79, 82, 86$. Median (12 values): $\frac{67 + 71}{2} = 69$. Group Y has the higher median.}
\end{practiceQuestion}

\begin{practiceQuestion}{m-8-3-q25}{gsa}
\begin{questionText}
The two circle graphs below compare how two families spend their monthly income.

\begin{center}
\begin{tikzpicture}
  % Family A: Housing 35%=126°, Food 25%=90°, Transport 20%=72°, Other 20%=72°
  \begin{scope}
  \node[above] at (0,2.3) {\textbf{Family A}};
  \fill[funBlue!70] (0,0) -- (0:2) arc (0:126:2) -- cycle;
  \fill[funRed!70] (0,0) -- (126:2) arc (126:216:2) -- cycle;
  \fill[funGreen!70] (0,0) -- (216:2) arc (216:288:2) -- cycle;
  \fill[yellow!70] (0,0) -- (288:2) arc (288:360:2) -- cycle;
  \draw (0,0) circle (2);
  \draw (0,0) -- (0:2); \draw (0,0) -- (126:2); \draw (0,0) -- (216:2); \draw (0,0) -- (288:2);
  \node at (63:1.2) {\small $35\%$};
  \node at (63:0.6) {\tiny Housing};
  \node at (171:1.2) {\small $25\%$};
  \node at (171:0.6) {\tiny Food};
  \node at (252:1.2) {\small $20\%$};
  \node at (252:0.6) {\tiny Trans.};
  \node at (324:1.2) {\small $20\%$};
  \node at (324:0.6) {\tiny Other};
  \end{scope}
  % Family B: Housing 40%=144°, Food 20%=72°, Transport 15%=54°, Other 25%=90°
  \begin{scope}[xshift=6cm]
  \node[above] at (0,2.3) {\textbf{Family B}};
  \fill[funBlue!70] (0,0) -- (0:2) arc (0:144:2) -- cycle;
  \fill[funRed!70] (0,0) -- (144:2) arc (144:216:2) -- cycle;
  \fill[funGreen!70] (0,0) -- (216:2) arc (216:270:2) -- cycle;
  \fill[yellow!70] (0,0) -- (270:2) arc (270:360:2) -- cycle;
  \draw (0,0) circle (2);
  \draw (0,0) -- (0:2); \draw (0,0) -- (144:2); \draw (0,0) -- (216:2); \draw (0,0) -- (270:2);
  \node at (72:1.2) {\small $40\%$};
  \node at (72:0.6) {\tiny Housing};
  \node at (180:1.2) {\small $20\%$};
  \node at (180:0.6) {\tiny Food};
  \node at (243:1.2) {\small $15\%$};
  \node at (243:0.6) {\tiny Trans.};
  \node at (315:1.2) {\small $25\%$};
  \node at (315:0.6) {\tiny Other};
  \end{scope}
\end{tikzpicture}
\end{center}

Family A earns $\$4{,}000$ per month and Family B earns $\$5{,}000$ per month. Which family spends more on housing, and by how much?
\end{questionText}
\correctAnswer{Family B spends more. Family A housing: $0.35 \times \$4{,}000 = \$1{,}400$. Family B housing: $0.40 \times \$5{,}000 = \$2{,}000$. Difference: $\$600$.}
\explanation{Even though Family A devotes a lower percentage to housing, the actual dollar comparison depends on income. Family B: $\$2{,}000 > \$1{,}400$. Difference: $\$2{,}000 - \$1{,}400 = \$600$.}
\end{practiceQuestion}

\begin{practiceQuestion}{m-8-3-q26}{gsa}
\begin{questionText}
The histograms below compare the ages of members at two community centers.

\begin{center}
\begin{tikzpicture}
\begin{axis}[
    ybar, bar width=14pt,
    width=6.5cm, height=5cm,
    ymin=0, ymax=14,
    xlabel={\small Age}, ylabel={\small Frequency},
    xtick={15,25,35,45,55},
    xticklabels={$10$--$19$,$20$--$29$,$30$--$39$,$40$--$49$,$50$--$59$},
    x tick label style={rotate=40, anchor=east, font=\tiny},
    title={\small Center A},
    nodes near coords, every node near coord/.style={font=\tiny},
]
\addplot[fill=funBlue!60] coordinates {(15,3) (25,5) (35,12) (45,7) (55,3)};
\end{axis}
\end{tikzpicture}
\hspace{0.5cm}
\begin{tikzpicture}
\begin{axis}[
    ybar, bar width=14pt,
    width=6.5cm, height=5cm,
    ymin=0, ymax=14,
    xlabel={\small Age}, ylabel={\small Frequency},
    xtick={15,25,35,45,55},
    xticklabels={$10$--$19$,$20$--$29$,$30$--$39$,$40$--$49$,$50$--$59$},
    x tick label style={rotate=40, anchor=east, font=\tiny},
    title={\small Center B},
    nodes near coords, every node near coord/.style={font=\tiny},
]
\addplot[fill=funRed!60] coordinates {(15,8) (25,10) (35,6) (45,4) (55,2)};
\end{axis}
\end{tikzpicture}
\end{center}

Compare the two distributions. Which center has a younger membership on average? Which has more variability in ages?
\end{questionText}
\correctAnswer{Center B has a younger membership (most members in the $10$--$29$ range). Center A has more variability — its ages are more spread out across all intervals, while Center B is concentrated at the younger end.}
\explanation{Center A peak: $30$--$39$ (middle-aged). Center B peak: $20$--$29$ (younger). Center A has significant frequencies across all intervals, while Center B is skewed toward younger ages, showing less spread.}
\end{practiceQuestion}

\begin{practiceQuestion}{m-8-3-q27}{gsa}
\begin{questionText}
The box plots below show daily high temperatures ($°$F) for two cities during one month.

\begin{center}
\begin{tikzpicture}
  \draw[->] (0,0) -- (9.5,0);
  \foreach \x/\v in {0/50,1/55,2/60,3/65,4/70,5/75,6/80,7/85,8/90}
    {\draw (\x,0.1) -- (\x,-0.1) node[below]{\tiny \v};}

  % City A: min=55, Q1=62, med=70, Q3=78, max=88
  \node[left] at (-0.3,1) {\small City A};
  \draw[thick] (1,0.7) -- (1,1.3);
  \draw[thick] (1,1) -- (2.4,1);
  \draw[thick,fill=funBlue!30] (2.4,0.7) rectangle (5.6,1.3);
  \draw[thick,funRed] (4,0.7) -- (4,1.3);
  \draw[thick] (5.6,1) -- (7.6,1);
  \draw[thick] (7.6,0.7) -- (7.6,1.3);

  % City B: min=60, Q1=66, med=72, Q3=76, max=82
  \node[left] at (-0.3,2.2) {\small City B};
  \draw[thick] (2,1.9) -- (2,2.5);
  \draw[thick] (2,2.2) -- (3.2,2.2);
  \draw[thick,fill=funGreen!30] (3.2,1.9) rectangle (5.2,2.5);
  \draw[thick,funRed] (4.4,1.9) -- (4.4,2.5);
  \draw[thick] (5.2,2.2) -- (6.4,2.2);
  \draw[thick] (6.4,1.9) -- (6.4,2.5);
\end{tikzpicture}
\end{center}

Find the IQR for each city. Which city has more consistent temperatures? Express the difference in medians as a multiple of City B's IQR.
\end{questionText}
\correctAnswer{City A IQR: $78 - 62 = 16°$F. City B IQR: $76 - 66 = 10°$F. City B is more consistent (smaller IQR). Medians: $72 - 70 = 2°$F difference $= \frac{2}{10} = 0.2$ IQRs — a small, likely insignificant difference.}
\explanation{City A: Q1$=62$, Q3$=78$, IQR$=16$. City B: Q1$=66$, Q3$=76$, IQR$=10$. City B has less variability. Median difference: $72-70=2$, which is $\frac{2}{10}=0.2$ of City B's IQR — very small.}
\end{practiceQuestion}
